\documentclass[a4paper]{article}
\input{../../../../../newpreamble.tex}
\title{Math 210A: Homework 6}
\author{Lance Remigio}
\begin{document}

\begin{problem}
    Let \( G  \) and \( J  \) be groups and let \( \varphi : G \to H  \) be a group epimorphism (surjective homomorphism)     
    \begin{enumerate}
        \item[(a)] Prove that if \( J  \) is abelian, then any subgroup of \( G  \) that contains \( \text{ker}(\varphi) \) is a normal subgroup of \( G  \).
            \begin{proof}
Note that the first isomorphism theorem says that 
   \[  G / \text{ker}(\varphi)  \cong \varphi(G) \]
   Also, \( \varphi(G) = J  \) since \( \varphi  \) is surjective. So, \( G / \text{ker}(\varphi) \cong J  \) and hence, \( G / \text{ker}(\varphi) \) is abelian. Since \( J  \) is abelian, every subgroup of \( G / \text{ker}(\varphi) \) is normal. Hence, by the 4th iso theorem every subgroup of \( G  \) containing \( \text{ker}(\varphi) \) is also normal.
            \end{proof}
        \item[(b)] Suppose \( K \trianglelefteq J  \). Prove that there exists \( H \trianglelefteq G  \) such that \( \text{ker}(\varphi) \subseteq  H  \) and \( G / H  \) is isomorphic to \( J / K  \).
            \begin{proof}
            Suppose \( K \trianglelefteq J  \). We will prove that our candidate for \( H  \) will be \( \varphi^{-1}(K) \leq  G  \). Indeed, we can see that \( \varphi^{-1}(K) \) is a subgroup of \( G  \) from homework 2. Now, our goal is to show that \( \varphi^{-1} \trianglelefteq G  \). It suffices to show that for all \( g \in G  \), 
            \[  g \varphi^{-1}(K) g^{-1} \subseteq  \varphi^{-1}(K). \]
            To this end, let \( g \in G  \) and let \( x \in g \varphi^{-1}(K) g^{-1} \). Then we have
            \[  x = g z g^{-1} \ \ \text{for some \( z \in \varphi^{-1}(K) \)}. \]
            Since \( z \in \varphi^{-1}(K) \), we have \( \varphi(z) \in K \). Then 
            \begin{align*}
                \varphi(x) = \varphi(g z g^{-1}) &= \varphi(g) \varphi(z) \varphi(g^{-1}) \\
                                                 &= \varphi(g) \varphi(z) (\varphi(g))^{-1} \\
                                                 &= j' \varphi(z) (j')^{-1}.
            \end{align*}
            Note that because \( K \trianglelefteq J  \), it follows that for all \( j \in J  \), 
            \[  j K j^{-1} \subseteq  K. \]
            In particular, for \( j = j' \) and since \( \varphi(z) \in K  \), we have that \( j' \varphi(z) (j')^{-1} \in K  \). This tells us that \( \varphi(x) \in K  \) and so \( x \in \varphi^{-1}(K) \). Thus, we see that, indeed, \( \varphi^{-1}(K) \trianglelefteq G  \). From part (a), we can see that \( \text{ker}(\varphi) \subseteq  \varphi^{-1}(K)  \). Hence, we can see that our candidate must be \( H = \varphi^{-1}(K) \).
            Now, our goal is to show that \( G / H \cong J / K  \). We will do this through the First Isomorphism Theorem. Define the mapping 
            \begin{center}
            \( \psi : G \to J / K  \) by \( g \mapsto \varphi(g) K   \).
            \end{center}
            Indeed, we can see that 
            \begin{align*}
                \text{ker}(\psi) &= \{ g \in G : \psi(g) = {1}_{H} K  \} \\
                                 &= \{ g \in G : \varphi(g) K = {1}_{H}  K  \}  \\
                                 &= \{  g \in G : \psi(g) \in K \}  \\
                                 &= \varphi^{-1}(K).
            \end{align*}
            Hence, we have \( \text{ker}(\psi) = \varphi^{-1}(K) = H  \). Because \( \varphi  \) is a surjective map, it also follows from our definition of \( \psi  \) that \( \psi  \) must also be surjective. Indeed, for any \( h \in H  \), there exists an \( g \in G  \) such that \( \varphi(g) = h  \). Hence, we have  
            \[  h K = \varphi(g) K  = \psi(g) \]
            and so \( \psi(G) = J / K  \).
            Using the First Isomorphism Theorem, we can see that 
            \[  G / \text{ker}(\psi) \cong \psi(G) \iff G / H  \cong J / K   \]
            where \( H = \varphi^{-1}(K)  \).

            \end{proof}
    \end{enumerate}
\end{problem}

\begin{problem}
   Suppose that \( \pi : {\Z}_{30} \to G  \) is an onto homomorphism where \( G  \) is a group of order \( 5  \). Determine, with proof, the kernel of \( \varphi \). 
\end{problem}
\begin{proof}
Observe, by the First Isomorphism Theorem, (because \( \pi  \) is an onto homomorphism) that  
\[  {\Z}_{30} / \text{ker}(\pi) \cong \pi({\Z}_{30}) \iff {\Z}_{30} / \text{ker}(\pi) \cong G.     \]
This tells us that, since \( G  \) and \( \Z_{30} \) are finite groups, it follows that 
\[  | {\Z}_{30} / \text{ker}(\pi) |  = | G  |  \iff \frac{ | {\Z}_{30} |  }{ |  \text{ker}(\pi) |  } = | G  |.   \]
Since \( | G  |  = 5  \) and \( | {\Z}_{30} | = 30  \), we have from the above that \( | \text{ker}(\varphi) | = 6   \). This tells us that there are six elements of \( {\Z}_{30} \) that gets mapped to the identity \( {1}_{G} \). Since the divisors of \( 30  \) gives us unique cyclic subgroups of \( {\Z}_{30} \), all we need to do is find the subgroup of the same order as \( \text{ker}(\varphi) \). Indeed, this subgroup of order \( 6  \) must be the cyclic subgroup generated by \( \langle  \overline{5} \rangle \). That is, \( \text{ker}(\pi) = \langle  \overline{5} \rangle = \{  \overline{0}, \overline{5}, \overline{10}, \overline{15}, \overline{20}, \overline{25} \}   \).
\end{proof}

\begin{problem}
    Suppose that \( | G  |  = p^{2} q  \) where \( p  \) and \( q  \) are primes with \( p > q  \). Show that if \( G  \) has a subgroup \( H  \) of order \( p^{2} \), then \( H  \) is the unique subgroup of order \( p^{2} \) in \( G  \). Deduce that \( G / H  \) is cyclic.
\end{problem}
\begin{proof}

\end{proof}

\begin{problem}
    Suppose that \( H  \) and \( K  \) are subgroups of finite index in a possibly infinite group \( G  \) with \( | G : H  |  = m  \) and \( | G : K  |  = n  \). Prove that \( \text{lcm}(m,n) \leq | G : H \cap K  |  \leq m n  \). Deduce that \( | G : H \cap K  |  = mn  \) when \( m  \) and \( n  \) are relatively prime.
\end{problem}
\begin{proof}

\end{proof}

\begin{problem}
    Prove Theorem 16 and corollary 17 from section 3.3 of the textbook.
\end{problem}

\begin{theorem}[The First Isomorphism Theorem]
   If \( \varphi : G \to H  \) is a homomorphism of groups, then \( \text{ker}(\varphi) \trianglelefteq G  \) and \( G / \text{ker}(\varphi) \cong \varphi(G) \). 
\end{theorem}
\begin{proof}
Define \( \mu : G / \text{ker}(\varphi) \to \varphi(G) \) by \( \mu(gK) = \varphi(g)  \) where \( K = \text{ker}(\varphi) \). To see that \( K  \) is normal, let \( g \in G  \) and let \( y \in g K g^{-1} \). Then \( y = g k g^{-1} \) for some \( k \in \text{ker}(\varphi) \). Then \( \varphi(k) = {1}_{H} \) and so 
\begin{align*}
    \varphi(y) = \varphi(g k g^{-1}) &= \varphi(g) \varphi(k) \varphi(g^{-1}) \\   
                                     &= \varphi(g) \varphi(k) (\varphi(g))^{-1} \\
                                     &= \varphi(g) {1}_{H} (\varphi(g))^{-1} \\
                                     &= \varphi(g)(\varphi(g))^{-1} \\
                                     &= {1}_{H}.
\end{align*}
Hence, it follows that \( y \in \text{ker}(\varphi) \) and so, \( K = \text{ker}(\varphi) \) must be a normal subgroup.
In what follows, we will show that \( \mu  \) isomorphism. First, we want to make sure that \( \mu  \) is well-defined. To this end, suppose \( {g}_{1} K = {g}_{2} K   \) where \( {g}_{1}, {g}_{2} \in G  \). This is true if and only if \( {g}_{1}^{-1}{g}_{2} \in \text{ker}(\varphi) \). This implies that \( \varphi({g}_{1}^{-1}) \varphi({g}_{2}) = {1}_{G}  \) which is true if and only if   
\[  (\varphi({g}_{1}))^{-1} \varphi({g}_{2}) = {1}_{G} \iff \varphi({g}_{2}) = \varphi({g}_{1}) \iff \mu({g}_{2}K) = \mu({g}_{1}K).  \]
This tells us that \( \mu  \) is well-defined and also that \( \mu  \) is injective.
Next, we will show that \( \mu  \) is a homomorphism. Indeed, for all \( {g}_{1} K , {g}_{2}  K \in G / K  \), we have that (since \( K  \) is a normal subgroup) 
\begin{align*}
    \mu({g}_{1}K {g}_{2}K) &= \mu({g}_{1}{g}_{2}K) \\ 
                           &= \varphi({g}_{1}{g}_{2})    \\
                           &= \big( \varphi({g}_{1}) \varphi({g}_{2}) \big)   \\
                           &= \varphi({g}_{1})  \varphi({g}_{2}) \\
                           &= \mu({g}_{1}K) \mu({g}_{2}K).
\end{align*}
Lastly, we show that \( \mu   \) is surjective. It is clear from the definition that this is the case. Indeed, for any \( y \in \varphi(g)  \), it follows that \( \varphi(g) = y  \) for some \( g \in G  \) and so we have 
\[  \mu(gK) = \varphi(g) = h.  \]
Hence, we see from the above argument that \( \mu  \) must be an isomorphism and so we conclude that  
\[  G / \text{ker}(\varphi) \cong \varphi(G).  \]
\end{proof}

\begin{corollary}[17]
   Let \( \varphi : G \to H  \) be a homomorphism of groups. Then 
   \begin{enumerate}
       \item[(1)] \( \varphi \) is injective if and only if \( \text{ker}(\varphi) = 1  \).
        \item[(2)] \( | G : \text{ker}(\varphi) |  = | \varphi(G) | \).
   \end{enumerate}
\end{corollary}
\begin{proof}
    First, we show (1). Since \( \varphi : G \to H  \) is a homomorphism, it follows from the First Isomorphism Theorem that
    \begin{align*}
        G / \text{ker}(\varphi) \cong \varphi(G).
    \end{align*}
    If \( \text{ker}(\varphi) = \{ {1}_{G} \}  \), then 
    \[  G / \text{ker}(\varphi) = G / \{ {1}_{G} \}  \cong G.  \]
    This holds if and only if
    \[  G \cong \varphi(G). \]
    which is true if and only if \( \varphi   \) is injective.
    Next, we show (2). Indeed, we can see from the First Isomorphism Theorem that (since \( \varphi \) is a homomorphism)
\begin{align*}
    G / \text{ker}(\varphi) \cong \varphi(G) &\iff | G / \text{ker}(\varphi) | = | \varphi(G) |  \\
                                             &\iff | G : \text{ker}(\varphi) |  = | \varphi(G) |.
\end{align*}
\end{proof}



\begin{problem}
    Let \( G  \) be a group with a normal subgroup \( N  \). Prove that there is a bijection from the subgroups of \( G  \) containing \( N \) onto the subgroup of \( G / N  \).
\end{problem}
\begin{proof}

\end{proof}

\end{document}

